\documentclass[12pt,a4paper]{article}
\usepackage[utf8]{inputenc}
\usepackage[T5]{fontenc}
\usepackage{amsmath, amssymb}
\usepackage{graphicx}
\usepackage{ifthen}

\newboolean{showsolutions}
\setboolean{showsolutions}{true}

% --- ĐỊNH NGHĨA CÁC LỆNH ẢO CHO CÔNG CỤ LATEX2JSON CỦA WEBSITE ---
\newcommand{\doctitle}[1]{\begin{center}\Large\textbf{#1}\end{center}}
\newcommand{\docdesc}[1]{\textbf{Mô tả:} #1}
\newcommand{\docgrade}[1]{\textbf{Lớp:} #1}
\newcommand{\docstatus}[1]{\textbf{Trạng thái:} #1}
\newcommand{\doctype}[1]{\textbf{Loại:} #1}
\newcommand{\doctopics}[1]{\textbf{Chủ đề:} #1}

\newenvironment{textblock}{\vspace{1em}}{}

\newenvironment{lesson}[2]{
	\vspace{1em}\noindent\textbf{\Large #1}\\
	\textit{#2}\vspace{0.5em}
}{}

\newcommand{\image}[3]{
	\begin{center}
		\includegraphics[width=#1\textwidth]{#3} \\
		\textit{#2}
	\end{center}
}

\newenvironment{quiz}[1]{
	\vspace{1em}\noindent\textbf{\Large Kiểm tra: #1}
}{}

\newenvironment{mcq}[4]{
	\vspace{1em}\noindent\textbf{Câu hỏi:} #1 \\
	\ifthenelse{\boolean{showsolutions}}{
		\textit{(Đáp án đúng: #2, Điểm: #3) - Giải thích: #4}
	}{}
	\begin{itemize}
	}{
	\end{itemize}
}
\newcommand{\option}[1]{\item #1}

\newenvironment{truefalse}[3]{
	\vspace{1em}\noindent\textbf{Đúng/Sai:} #1 \\
	\ifthenelse{\boolean{showsolutions}}{
		\textit{(Điểm: #2) - Giải thích: #3}
	}{}
	\begin{itemize}
	}{
	\end{itemize}
}
\newcommand{\statement}[2]{\item [\textbf{#1}] #2}

\newcommand{\shortanswer}[4]{
    \vspace{1em}\noindent\textbf{Trả lời ngắn (Điểm: #3):}

    \noindent #1

	\ifthenelse{\boolean{showsolutions}}{
		\vspace{0.5em}
		\noindent\textit{Đáp án: #2 - Giải thích:}
		\noindent #4
	}{}
}

\newcommand{\essay}[3]{
    \vspace{1em}\noindent\textbf{Tự luận (Điểm: #2):}
    
    \noindent #1
    
	\ifthenelse{\boolean{showsolutions}}{
		\vspace{0.5em}
		\noindent\textbf{Gợi ý / Đáp án:}
		\noindent #3
	}{}
}

\begin{document}

\doctitle{ĐỀ LUYỆN TẬP SỐ 02: PHƯƠNG TRÌNH MẶT PHẲNG}
\docdesc{Bộ đề trắc nghiệm rèn luyện phương trình mặt phẳng trong không gian Oxyz - Đề số 02}
\docgrade{12}
\docstatus{draft}
\doctype{test}
\doctopics{Hình học không gian, Phương trình mặt phẳng}

% =========================================================================
% PHẦN 1. CÂU TRẮC NGHIỆM NHIỀU PHƯƠNG ÁN LỰA CHỌN
% =========================================================================

\begin{quiz}{Phần 1. Câu trắc nghiệm nhiều phương án lựa chọn}

\begin{mcq}{Trong không gian $Oxyz$, cho mặt phẳng $(\alpha): 2x + y - 3z + 8 = 0$. Mặt phẳng nào sau đây vuông góc với mặt phẳng $(\alpha)$?}{1}{1}{Chọn B.}
	\option{$x - 3y + 3z - 7 = 0$.}
	\option{$3x - 3y + z - 7 = 0$.}
	\option{$x + 2y - z - 8 = 0$.}
	\option{$x - 2y + z + 8 = 0$.}
\end{mcq}

\begin{mcq}{Trong không gian $Oxyz$, tìm phương trình mặt phẳng $(\alpha)$ cắt ba trục $Ox, Oy, Oz$ lần lượt tại ba điểm $A(-3; 0; 0), B(0; 4; 0), C(0; 0; -2)$.}{2}{1}{Chọn C.}
	\option{$4x - 3y + 6z - 12 = 0$.}
	\option{$4x + 3y - 6z + 12 = 0$.}
	\option{$4x - 3y + 6z + 12 = 0$.}
	\option{$4x + 3y + 6z + 12 = 0$.}
\end{mcq}

\begin{mcq}{Trong không gian $Oxyz$, mặt phẳng $(P)$ đi qua điểm $G(1; 1; 1)$ và vuông góc với đường thẳng $OG$ có phương trình là:}{0}{1}{Chọn A.}
	\option{$x + y + z - 3 = 0$.}
	\option{$x - y + z = 0$.}
	\option{$x + y - z - 3 = 0$.}
	\option{$x + y + z = 0$.}
\end{mcq}

\begin{mcq}{Trong không gian $Oxyz$, mặt phẳng $(P)$ đi qua $M(-1; 2; 4)$ và chứa trục $Oy$ có phương trình là:}{1}{1}{Chọn B.}
	\option{$4x - z = 0$.}
	\option{$4x + z = 0$.}
	\option{$x - 4z = 0$.}
	\option{$x + 4z = 0$.}
\end{mcq}

\begin{mcq}{Trong không gian $Oxyz$, khoảng cách từ điểm $A(1; -2; 3)$ đến $(P): x + 3y - 4z + 9 = 0$ bằng:}{3}{1}{Chọn D.}
	\option{$\frac{\sqrt{26}}{13}$.}
	\option{$\sqrt{8}$.}
	\option{$\frac{17}{\sqrt{26}}$.}
	\option{$\frac{4\sqrt{26}}{13}$.}
\end{mcq}

\begin{mcq}{Trong không gian $Oxyz$, cho hai mặt phẳng $(\alpha): x - 2y - 2z + 4 = 0$ và $(\beta): -x + 2y + 2z - 7 = 0$. Tính khoảng cách giữa hai mặt phẳng $(\alpha)$ và $(\beta)$.}{3}{1}{Chọn D.}
	\option{$3$.}
	\option{$-1$.}
	\option{$0$.}
	\option{$1$.}
\end{mcq}

\begin{mcq}{Trong không gian $Oxyz$, mặt phẳng nào sau đây song song với $(P): 2x - 4y + 6z - 5 = 0$?}{3}{1}{Chọn D.}
	\option{$(\alpha): 2x - 4y - 6z - 5 = 0$.}
	\option{$(\beta): x - 2y + 3z - \frac{5}{2} = 0$.}
	\option{$(\chi): x + 2y + 3z - \frac{5}{2} = 0$.}
	\option{$(\gamma): x - 2y + 3z - 5 = 0$.}
\end{mcq}

\begin{mcq}{Trong không gian $Oxyz$, cho hai mặt phẳng $(P): 2x - 5y + 4z + 3 = 0$ và $(Q): 3x - 2y - 4z - 21 = 0$. Vị trí tương đối của $(P)$ và $(Q)$ là:}{3}{1}{Chọn D.}
	\option{Song song với nhau.}
	\option{Trùng nhau.}
	\option{Cắt nhưng không vuông góc nhau.}
	\option{Vuông góc nhau.}
\end{mcq}

\end{quiz}

% =========================================================================
% PHẦN 2. CÂU TRẮC NGHIỆM ĐÚNG SAI
% =========================================================================

\begin{quiz}{Phần 2. Câu trắc nghiệm đúng sai}

\begin{truefalse}{Trong không gian $Oxyz$, cho điểm $A(1; 2; -1)$ và hai mặt phẳng $(P), (Q)$ lần lượt có phương trình là $2x + y - 2z = 0$ và $-8x - 4y + 8z + 6 = 0$. Xét tính đúng sai của các khẳng định sau:}{1}{
- Thay $A(1; 2; -1)$ vào $(P)$ được $2(1) + 2 - 2(-1) = 6 \neq 0 \Rightarrow A \notin (P)$. Thay vào $(Q)$ được $-8(1) - 4(2) + 8(-1) + 6 = -18 \neq 0 \Rightarrow A \notin (Q)$. Khẳng định a đúng.
- $(Q) \Leftrightarrow 2x + y - 2z - \frac{3}{2} = 0$ song song với $(P)$. Khẳng định b đúng.
- $d((P), (Q)) = \frac{|0 - (-3/2)|}{\sqrt{2^2 + 1^2 + (-2)^2}} = \frac{1}{2} = 0{,}5$ không phải số nguyên. Khẳng định c sai.
- $d(A, (P)) = 2, d(A, (Q)) = \frac{3}{2} \Rightarrow d(A, (P)) = \frac{4}{3} d(A, (Q))$. Khẳng định d đúng.
}
	\statement{true}{Điểm $A$ không thuộc cả hai mặt phẳng $(P)$ và $(Q)$.}
	\statement{true}{Hai mặt phẳng $(P)$ và $(Q)$ song song với nhau.}
	\statement{false}{Khoảng cách giữa hai mặt phẳng $(P)$ và $(Q)$ là một số nguyên dương.}
	\statement{true}{$d(A, (P)) = \frac{4}{3} d(A, (Q))$.}
\end{truefalse}

\begin{truefalse}{Trong không gian $Oxyz$, cho ba điểm $A(1; 0; 0), B(0; 0; 2), C(0; -3; 0)$. Xét tính đúng sai của các khẳng định sau:}{1}{
- Phương trình đoạn chắn $(ABC): \frac{x}{1} + \frac{y}{-3} + \frac{z}{2} = 1 \Leftrightarrow 6x - 2y + 3z - 6 = 0$. VTPT là $(6; -2; 3)$ không cùng phương với $(1; -3; 2)$. Khẳng định a sai.
- Thay $M(1; 3; 2)$ vào $(ABC)$ được $6(1) - 2(3) + 3(2) - 6 = 0 \Rightarrow M \in (ABC)$. Khẳng định b đúng.
- $d(O, (ABC)) = \frac{|-6|}{\sqrt{6^2 + (-2)^2 + 3^2}} = \frac{6}{7}$. Khẳng định c đúng.
- Mặt phẳng qua $A(1; 0; 0)$ vuông góc với $BC$ nhận $\vec{BC} = (0; -3; -2)$ làm VTPT có phương trình $0(x - 1) - 3(y - 0) - 2(z - 0) = 0 \Leftrightarrow 3y + 2z = 0$. Khẳng định d đúng.
}
	\statement{false}{Một vectơ pháp tuyến của mặt phẳng $(ABC)$ là $\vec{n} = (1; -3; 2)$.}
	\statement{true}{Mặt phẳng $(ABC)$ đi qua điểm $M(1; 3; 2)$.}
	\statement{true}{$d(O, (ABC)) = \frac{6}{7}$.}
	\statement{true}{Mặt phẳng đi qua điểm $A$ và vuông góc với $BC$ có phương trình là $3y + 2z = 0$.}
\end{truefalse}

\end{quiz}

% =========================================================================
% PHẦN 3. CÂU TRẮC NGHIỆM TRẢ LỜI NGẮN
% =========================================================================

\begin{quiz}{Phần 3. Câu trắc nghiệm trả lời ngắn}

\shortanswer{Trong không gian $Oxyz$, cho hình chóp $S.ABCD$ có đáy là hình chữ nhật và các điểm $A(0; 0; 0), B(1; 0; 0), D(0; 2; 0), S(0; 0; 3)$. Khoảng cách từ trọng tâm $G$ của tam giác $SAB$ đến mặt phẳng $(SBD)$ bằng bao nhiêu? (Làm tròn kết quả đến hàng phần chục).}{0,3}{1}{Trọng tâm $G\left(\frac{1}{3}; 0; 1\right)$. Mặt phẳng $(SBD): 6x + 3y + 2z - 6 = 0$. Khoảng cách $d(G, (SBD)) = \frac{|2 + 0 + 2 - 6|}{\sqrt{6^2 + 3^2 + 2^2}} = \frac{2}{7} \approx 0{,}3$. Đáp án: 0,3.}

\shortanswer{Trong không gian $Oxyz$, cho điểm $M(-1; 0; 3)$. Hỏi có bao nhiêu mặt phẳng $(P)$ đi qua điểm $M$ và cắt các trục $Ox, Oy, Oz$ lần lượt tại $A, B, C$ sao cho $3OA = 2OB = OC \neq 0$?}{2}{1}{Ta có $|a| = \frac{k}{3}, |b| = \frac{k}{2}, |c| = k$. Thay $M(-1; 0; 3)$ vào $(P): \frac{x}{a} + \frac{y}{b} + \frac{z}{c} = 1 \Rightarrow -\frac{1}{a} + \frac{3}{c} = 1$. Chỉ có $a = -\frac{k}{3}, c = k \Rightarrow k = 6 \Rightarrow a = -2, c = 6$ và $b = \pm 3$. Có 2 mặt phẳng thỏa mãn. Đáp án: 2.}

\shortanswer{Khi gắn hệ trục tọa độ $Oxyz$ (đơn vị trên mỗi trục tọa độ là decimét) vào một ngôi nhà 1 tầng, người ta thấy rằng mặt trên và mặt dưới của mái nhà thuộc các mặt phẳng vuông góc với trục $Oz$. Biết rằng các vị trí $A(3; 4; 33), D(9; 8; 35)$ lần lượt thuộc mặt dưới, mặt trên của mái nhà. Độ dày của mái nhà được tính bằng khoảng cách giữa mặt trên và mặt dưới của mái nhà đó. Hãy cho biết độ dày của mái nhà đó là bao nhiêu decimét?}{2}{1}{Mặt dưới $(\alpha): z = 33$, mặt trên $(\beta): z = 35$. Độ dày của mái nhà là $d((\alpha), (\beta)) = |35 - 33| = 2\text{ dm}$. Đáp án: 2.}

\end{quiz}

\end{document}